%
% Hi,
% after reading through the Introduction of the thesis.pdf welcome to its source.
% This is the main file for your thesis.
% It mainly loads other things and gets the settings right.
%
% Before going on a word about editing latex:
% - encoding: use utf8 encoding, no excuse. If your editor can't use another one.
% - revision control: you want it, it saves your life, no matter which one (i suggest git / mercurial)
% 
% Especially important with multiple authors:
% - revision control: most of them use line based diffs
%                     In text (not code) i suggest to write one sentence per line (so hit ENTER after each sentence).
%                     Use dynamic word wrapping in your editor so you don't have to scroll all the time.
%                     This eases merging and does not show a lot of diffs just due to reformatting (as with static word wrapping).
%
% I wrote my thesis in English and so should you if you ever want to reuse or publish your results (or parts of it).
% Even in the English thesis you might want to include a German abstract and your "Erklärung" has to be in German anyhow, which will be mentioned in the comments below.
% If you still want to write it in German: no problem there are just a few things to change in this file which will also be mentioned in the comments.
%
% So how to start?
% Go through this file, it includes quite a lot of packages which you might want to use or not.
% Below all the packages you'll find the \begin{document} line.
% Underneath it insert your title, subtitle, name, supervisors, etc.
% Then have a look into the files it includes (starting with abstract, etc.
% 
% A "how to find documentation" tip:
% Usually you can quickly find help for some package by typing:
% texdoc packagename
% into your terminal.
% try: texdoc scrbook
% 
% 
% 
% Jörn
%
% PS:
% read this: ftp://ftp.dante.de/tex-archive/info/l2tabu/german/l2tabu.pdf

\RequirePackage{fix-cm} % see: http://tug.org/texmf-dist/doc/latex/base/fixltx2e.pdf
\documentclass[a4paper, twoside, BCOR=10mm]{scrbook} % see docs!

%\usepackage[ngerman]{babel} % if your whole thesis is German uncomment
%%\usepackage{ngerman} % actually deprecated by babel package with ngerman option (line above), do not uncomment
\usepackage[ngerman, USenglish]{babel} % main language comes last!

\usepackage[T1]{fontenc}
\usepackage{lmodern} % modern fonts
%\usepackage{textcomp}
\usepackage{fixltx2e} % some other fixes

\usepackage[utf8]{inputenc} % This is the only plausible encoding. if you can't read äöüß now, CHANGE YOUR EDITOR's ENCODING, DO NOT CHANGE THIS LINE (only exception: on some systems it's utf8x, so in case you get an error try that)

\usepackage{scrhack} % some fixes for the scr-classes

%\usepackage[german]{varioref} % References with page numbers
\usepackage{varioref} % References with page numbers (use \vref instead of \ref)

% \usepackage[printonlyused,withpage]{acronym}
\usepackage{acronym} % acronyms, first occurrence will be expanded, all further just acronym + acronym table in appendix

%\usepackage{setspace} % (usually not a good idea)
%\doublespacing
%\onehalfspacing

% für das Euro symbol (\euro{}) oder Beträge: \EUR{47,11}
\usepackage{eurosym}

% \usepackage{floatflt} %floating figures (Figures/Tables, surrounded by text)
\usepackage{subfigure} % multiple figures (and references in one figure)

\usepackage{graphicx} % to import graphics


%\usepackage{qtree} % draw trees in latex

% \usepackage[active]{srcltx} % deprecated: ForwardDVI find current position in .dvi
% (ForwardPDF) current Editors/Distributions use pdflatex with option -interaction=nonstopmode -synctex=1 '%source'

% in case urls are included
\usepackage{url}
%\usepackage{hyperref} % clickable links
\usepackage[colorlinks=true,
            linkcolor=black,
            anchorcolor=black,
            citecolor=black,
            filecolor=black,
            menucolor=black,
            runcolor=black,
            urlcolor=black]{hyperref} % without all those nasty colored boxes in pdf viewer



% math stuff

\usepackage{amsmath}
%\usepackage{mathtools}

\usepackage{amsfonts}
\usepackage{amssymb}
\usepackage{amsthm}
%\usepackage{eulervm}

%\DeclareMathSymbol{\R}{\mathalpha}{AMSb}{"52} % nasty like this, use the following lines
\newcommand{\R}{\mathbb{R}}
\newcommand{\F}{\mathbb{F}}
\newcommand{\colonequals}{\mathrel{\mathop:}=}


% coding stuff

\usepackage{color} % used by listings for syntax highlighting
\definecolor{darkblue}{rgb}{0,0,.6}
\definecolor{darkred}{rgb}{.6,0,0}
\definecolor{darkgreen}{rgb}{0,.6,0}
\definecolor{red}{rgb}{.98,0,0}

\usepackage{verbatim}
\usepackage{listings} % for code snippets, see docs
\renewcommand{\lstlistlistingname}{List of Listings}
\lstset{
	commentstyle=\itshape\color{darkgreen},
	keywordstyle=\bfseries\color{darkblue},
	stringstyle=\color{darkred},
	showspaces=false,
	showtabs=false,
	showstringspaces=false,
% 	xleftmargin=1cm,
	columns=fixed,
	frame=single,
	basicstyle=\ttfamily\small,
	language=Python,
	breaklines=true,
	tabsize=3,
% 	numbers=left,
% 	numberfirstline=true,
%	firstnumber=0,
% 	stepnumber=5,
	extendedchars=true,
	captionpos=b,
}

% if you want line numbers like 1,5,10,15,... here's a hack (thanks to kaste in #latex on freenode):
% \lstset{numbers=left, numberstyle=\tiny, stepnumber=5, numbersep=5pt, numberfirstline=true}
% \begin{document}
% \hbox to 0pt{\begin{lstlisting}[firstnumber=0, name=test]
%         
% \end{lstlisting}}
% \begin{lstlisting}[name=test]
% for i:=maxint to 0 do
% begin
%      { do nothing }
% end;
% more
% crap
% for i:=maxint to 0 do
% begin
%      { do nothing }
% end;
% more
% crap\end{lstlisting}
% \end{document}


\usepackage{paralist} % provides inline lists (so an enumeration inline in a paragraph with \begin{inparaenum}[(i)]...\end{...}

\usepackage{rotating} % e.g., if you want a landscape figure

% array stuff
\usepackage{array}
\usepackage{longtable} % multi-page tables
\usepackage{tabularx}
\usepackage{tabulary}
\usepackage{multirow} % cells spanning multiple rows
\usepackage{multicol} % cells spanning multiple columns

\KOMAoption{toc}{bib,idx} %,listof} % what is part of the Table Of Contents? (needs to be defined after listings package but before natbit!)

%\usepackage{bibgerm}
%\usepackage{authordate1-4}
\usepackage[square]{natbib} % use the natbib library package, see docs!
%\bibpunct[,]{[}{]}{;}{a}{,}{,} %[notesep],open,close,citesep,style,aysep, yysep
\renewcommand{\cite}{\citep} % by default use natbib's citep wherever we write \cite{...}


%\usepackage{lineno} % sometimes useful for reviews
%%\linenumbers
%\pagewiselinenumbers
%\modulolinenumbers[5]
%\switchlinenumbers % Doppelseitig außen (\switchlinenumbers* % Doppelseitig innen)


% word index -> page numbers for appendix.
\makeindex
\usepackage{makeidx}


% sometimes useful to just get some random text on a page
%\usepackage{blindtext} % for sampletext (notice that you need to use \usepackage[english]{babel} with this)





% NICE TO KNOW:
% Befehl: (a, b are words)
% a.\ b    A non-expandable space (e.g., c.f.~here) (usually you should use this after abbreviation "."s
% a~b      A non-breakable space between a and b (e.g., Figure~\ref{fig:blabla})
% a\,b     Small space (e.g., with Web\,2.0)
% a\-b     Hyphenation: only at this point
% a"-b     Hyphenation: additionaly at this point
% a""b     Hyphenation: at this point without hyphen
% a-b      Hyphenation: always print a hyphen, permitting no other hyphenations
% a"~b     Hyphenation: hyphen here, but no hyphenation allowed at it
% a"=b     Hyphenation: always print a hyphen here but allow other hyphenations
% ``a''    English caps
% "`a"'    German caps

% howto use varioref:
% \vref{} instead of \ref{}

% howto use natbib:
% \citet{jon90}               ⇒ Jones et al. (1990)
% \citet[chap.~2]{jon90}      ⇒ Jones et al. (1990, chap. 2)
% \citep{jon90}               ⇒ (Jones et al., 1990)
% \citep[chap.~2]{jon90}      ⇒ (Jones et al., 1990, chap. 2)
% \citep[see][]{jon90}        ⇒ (see Jones et al., 1990)
% \citep[see][chap.~2]{jon90} ⇒ (see Jones et al., 1990, chap. 2)
% \citet*{jon90}              ⇒ Jones, Baker, and Williams (1990)
% \citep*{jon90}              ⇒ (Jones, Baker, and Williams, 1990)
% \citealt{jon90}             ⇒ Jones et al. 1990
% \citealt*{jon90}            ⇒ Jones, Baker, and Williams 1990
% \citealp{jon90}             ⇒ Jones et al., 1990
% \citealp*{jon90}            ⇒ Jones, Baker, and Williams, 1990
% \citealp{jon90,jam91}       ⇒ Jones et al., 1990; James et al., 1991
% \citealp[pg.~32]{jon90}     ⇒ Jones et al., 1990, pg. 32
% \citenum{jon90}             ⇒ 11
% \citetext{priv.\ comm.}     ⇒ (priv. comm.)
% \citeauthor{jon90}          ⇒ Jones et al.
% \citeauthor*{jon90}         ⇒ Jones, Baker, and Williams
% \citeyear{jon90}            ⇒ 1990
% \citeyearpar{jon90}         ⇒ (1990)
% \citet{dRob98}              ⇒ della Robbia (1998)
% \Citet{dRob98}              ⇒ Della Robbia (1998)
% \Citep{dRob98}              ⇒ (Della Robbia, 1998)
% \Citealt{dRob98}            ⇒ Della Robbia 1998
% \Citealp{dRob98}            ⇒ Della Robbia, 1998
% \Citeauthor{dRob98}         ⇒ Della Robbia



%\hyphenation{My-word-which-shall-always-only-be-seperated-here and-another-word and-a-third-one and-so-on} % if you spot mis-hyphenations in your whole document, insert them here.



\begin{document}
%Userdefs (change these, as they will be inserted in various places e.g., title)
\newcommand{\ThesisType}{Master Thesis} % or "Bachelor Thesis"
\newcommand{\TheTitle}{Example Title}
\newcommand{\TheSubTitle}{Example Subtitle}
\newcommand{\TheAuthor}{Example Author}
\newcommand{\Email}{\url{exampleEmail@cs.uni-kl.de}}
\newcommand{\betreuerA}{Prof. Dr. My Example Prof}
\newcommand{\betreuerB}{Dr. My Example Dr}
\newcommand{\abgabedatum}{April\ 1st, 2011} % English date (same as German for German thesis)
\newcommand{\abgabedatumDE}{1. April 2011} % German date (used in "erklaerung.tex")


% don't change these or know what you do
\subject{\ThesisType}
\title{\TheTitle}
\subtitle{\TheSubTitle}
\author{\TheAuthor}
\date{\abgabedatum}

\hypersetup{pdftitle={\TheTitle: \TheSubTitle},
            pdfauthor={\TheAuthor},
            pdfcreator={\TheAuthor},
            pdfsubject={\ThesisType} }







%%%%%%%%%%%%%%%%%%%%%%%%%%%%%%%%%%%%%%%%%%%%%%%%%%%%%%%%%%%%%%%%%%%%%%%%%%%%%%
%                              Usercommands
%%%%%%%%%%%%%%%%%%%%%%%%%%%%%%%%%%%%%%%%%%%%%%%%%%%%%%%%%%%%%%%%%%%%%%%%%%%%%%



%Insert a footnote with reference to label including page numbers.
\newcommand{\vgl}[1]{\footnote{sf. Chapter \vref{#1}.}}



% editing & structuring
\newcommand{\redline}[1]{\ignorespaces\begin{scriptsize}\begin{itemize}\setlength{\parskip}{0pt}%
	\setlength{\parsep}{0pt}\color{darkred}#1\end{itemize}\end{scriptsize}\ignorespaces}
% \newcommand{\redline}[1]{}

% good for writing a quick outline into your thesis... for example like this:
% \section{blabla}
% \redline{\item first point i want to mention
%          \item second point, etc.}
% either comment out as soon as you have written the content of the section or just redefine redline to nothing as above (i like to comment things out more).



\newcommand{\introduce}[1]{\emph{#1}\index{#1}} % first time you mention a specific concept (appears in index)
\newcommand{\reuse}[1]{\emph{#1}\index{#1}} % if you want to add something to its definition or the like (reappears in index)

\newcommand{\lastAccessed}{(last accessed: Jan.\ 23rd, 2011)} % e.g., for URIs

% math formula:
\newcommand{\lO}{\mathcal{O}} % Landau O

% specifics (feel free to put a lot more here)
\newcommand{\curie}[1]{\url{#1}} 
\newcommand{\iri}[1]{<\url{#1}>}
\newcommand{\iripart}[1]{\url{#1}}



\frontmatter % switch scrbook parts

%\thispagestyle{empty}

% you don't need to change anything in this file, as it is all taken from the definitions in thesis.tex
\begin{titlepage}

\parbox[t]{12cm}{
	\parbox[c]{5cm}{\includegraphics[width=5cm]{pics/Logo-TU-KL-text}}
	
	\vspace{1ex}
	\hspace{1.39cm}
	\parbox[t]{8.5cm}{\textsf{\textit{University of Kaiserslautern}\\[0.4cm]
		\textbf{Fachbereich Informatik}\\
		\textit{Department of Computer Science}}}
}

\begin{center}
	\vspace{3cm}
	
	\large
	\textbf{\ThesisType}
	
	
	\vspace{3em}
	
	\LARGE
	\textbf{\TheTitle:\\
	\Large \TheSubTitle}
	
	\large
	
	\vspace{5em}
	
	\textbf{\TheAuthor}\\[1ex]
	{\small{\Email}}
	
	\vfill
	\abgabedatum
\end{center}

\vfill

\begin{center}
	\renewcommand*{\and}{ \\ &}
	\begin{tabular}{ll}
		Supervisors: & \betreuerA \\
		             & \betreuerB
		
	\end{tabular}
\end{center}
%\newpage

\end{titlepage}


% Abstracts part:
% If you want English only don't touch.
% If you want _additional_ German abstract: Uncomment everything below.
% If you just want a German abstract: 
% if you want abstracts on left and right page, look up what is commented out here
%\cleardoubleevenemptypage
\begingroup
	\chapter{Abstract}
	This is a skeleton of my master thesis.
It includes a quite common structure and some details on how to use it in the Introduction.

On this page you can see an abstract.
It shall be short (max. 1 page) and give an uninterested reader a clue what this whole work is about.
Why do you think your work is needed?
What is your main contribution?

Again: make it short!

	%\renewcommand*{\cleardoublepage}{\clearpage}
	%\clearpage
% 	\begin{otherlanguage}{ngerman} % uncomment this line if you need an additional
	                               % German abstract in an English thesis
% 	\chapter*{Zusammenfassung} % (also) uncomment if you need a German abstract
% 	Hier käme die Zusammenfassung der Arbeit auf deutsch hin.

        % ...
% 	\end{otherlanguage}
\endgroup
%\cleardoubleevenemptypage

\setcounter{tocdepth}{1} % so only chapters and sections (best imho)
\tableofcontents

% lists on one page
\chapter{Lists}
%\markboth{Lists}{Lists}
\begingroup
	\renewcommand*{\chapter}{\section}
	\renewcommand*{\addchap}{\addsec}
	%\renewcommand*{\markboth}[1]{\markright}
	\renewcommand*{\listfigurename}{Figures}
	\listoffigures
	\renewcommand*{\listtablename}{Tables}
	\listoftables
	\renewcommand{\lstlistlistingname}{Listings}
	\lstlistoflistings
\endgroup

\cleardoubleevenemptypage
\mainmatter % switch scrbook parts, numbered: 1., 2., ...
\acresetall


\chapter{Introduction}
	Start with a motivation.
	What is the problem?
	What is it part of?
	What is the exact problem you want to solve in your work?
	What is the benefit of solving it?
	How do you solve it (sketch only)?

% 	\section{Research Question}
% 		%  1. a concise statement of the question that your thesis tackles 
% 		%  2. justification, by direct reference to section 3 that your question is previously unanswered
% 		%  3. discussion of why it is worthwhile to answer this question.
	
	\section{About This Skeleton}
		This skeleton tries to guide you in finding a good structure for your thesis and help you to represent it in \LaTeX.
		
		You should be aware of the fact that each field of computer science comes with its own ``default structure'' for a thesis.
		The default structure of your thesis has a very central purpose for your reader: quickly finding the relevant parts and getting an overview.
		I heard arguments for not following this default structure, because it is boring.
		I don't care, it's not a novel or a story you're writing, you're focus lies on quickly exposing information to your reader.
		
		If I was to read (not necessarily grade) your thesis I would read it like this (pretty much like a paper):
		\begin{description}
			\item[Abstract] (What is it about? What are the central ideas? Do I care/need this?)
			\item[Table of contents] (Short glimpse where to expect what)
			\item[Intro] (Search for main contribution and ``why?'')
			\item skip over the thesis in large steps while going to the conclusion (Just for orientation)
			\item[Conclusion] (Does it work, problems solved, future work, discussion)
			\item[References] (Do you cite key papers in the area?)
			\item[Evaluation] (Is the evaluation sound? does it really work?)
			\item[Related work] (What is related? how does your approach differ?)
			\item[Main parts] (Details on how you solved things)
			\item[Foundations] (Only if referred to from main parts and things are unknown to me)
		\end{description}
		In any of these steps i would stop if i lost my interest.
		This probably doesn't count for someone grading your thesis as (s)he has to go on anyhow, but probably these still are the first steps (and the last two are reading the work completely).
		
		As mentioned before every subfield of computer science seems to have its own structure.
		Math works look quite different from software engineering works and implementation works look different as well.
		Still the main structure can often be found maybe with different names (e.g., evaluation might be called complexity analysis...).
		
		So let yourself be guided by this skeleton, try to find something for each of the chapters and possibly rename them according to your needs.
		And last but not least: 
		Make sure to check back with your supervisors in an early stage e.g., after a first ``extended outline'' (mainly chapters, sections, subsections and \verb|\redline|s), before wasting a lot of time for nothing.
	
	\section{Structure of this Thesis}\label{sec:structure} 
% 		\redline{\item Figure~\ref{fig:ProcessOverview}
% 					\item Input: Data Acquisition and Preprocessing
% 					\item The Game: Getting player preferences
% 					\item Output: Generate a Ranking from acquired preferences}
% 		\begin{figure}
% 			\centering
% 			\includegraphics[width=\textwidth]{./pics/processOverview}
% 			\caption{Data Flow Process Overview}
% 			\label{fig:ProcessOverview}
% 		\end{figure}
		
		Briefly explain the structure of the thesis.
		If you have a process overview, include a picture showing the whole thing (better than a thousand words).
		Explain which part of the picture is explained where leading do your evaluation and conclusion.
		



\chapter{Foundations}\label{chap:Foundations}
	A chapter for (mathematical) foundations and background knowledge your audience will need.
	
	\section{What Belongs Here?}
	While writing this section usually teaches you a lot (as you have to explain the concepts you are using), keep in mind that this is \underline{not} your actual work.
	This chapter should be as short as possible, brushing up on things which are essential for your work.
	
	\section{What Does Not}
	Do not forget that your reader usually is a computer scientist, so please do not re-explain too simple things here.
	

\chapter{Related Work}\label{chap:relatedWork}
	This is an important chapter and it is different from Chapter~\ref{chap:Foundations}.
	
	While you focus on what things a reader should know to understand your work in the previous chapter, this chapter focuses on previous work which is (very) close to the problem you want to solve.
	
	Maybe previous authors nearly solved the problem you want to solve or maybe they solved in a different way.
	Mention such works here (e.g., \cite{Bush1945}) and tell the reader what makes your work special.
	Show the \underline{differences}, how is your approach better or worse?
	
	
	\section{Field A}
		If you're combining works from several fields, why not structure them into Sections?
	
	\section{Field B}
		And another field.
	

\chapter{Mainpart}
	This is just a dummy chapter for all the main chapters of your thesis, so the actual content of your work.
	
	Usually it's a poor choice to just have one main chapter and a huge number of sections.
	Also you don't name the Mainpart Mainpart, give these chapters speaking names.
	Remember Section~\ref{sec:structure}.
	 % possibly split up into several files
%\include{otherfiles}

\chapter{Evaluation}\label{chap:eval}
	Evaluate your approach.
	Compare it to others (Chapter~\ref{chap:relatedWork}), give us stats, etc.
	
	If your work is more theoretical you might want to rename this section to Analysis and for example include the complexity analysis in here.


\chapter{Conclusion}\label{chap:conclusion}
	Conclude your work.
	Summarize the results and findings, perhaps discuss the evaluation results here and say what is left to be done.

	\section{Summary}\label{sec:summary}
		
% 	\section{Discussion}\label{sec:discussion}
		
	\section{Open Questions and Future Work}
	

%\nocite{*} % if you want to include everything from your library even if never \cited (usually bad idea)

% some bibliographystyles you might want to give a try... i like mine better ;) (these just set the style for your library without including it)
% \bibliographystyle{plainnat} % plain (and a lot)
% \bibliographystyle{apalike} % a1 and a2; a1, a2, and a3 (lastname, f.)
%\bibliographystyle{apalike2} % a1 & a2; a1, a2, & a3 (lastname, f.)
% \bibliographystyle{newapa} %  a1 & a2; a1, a2, & a3 (lastname, f.)
% \bibliographystyle{chicago} 
% \bibliographystyle{named}

\bibliographystyle{joernstyle} % based on apalike2 and plainnat, prints URLs for misc, works with mendeley bibtex exports

% while writing you should use this
\bibliography{library} % actually includes the library (here)
% when your work is done, you should copy the thesis.bbl to
% a file called bibFrozen.tex and include that, with that your thesis lib won't
% change anymore, independently from possible changes to your library.bib or 
% bibtex versions or styles or different platforms... you get it.
%
% Maybe even edit bibFrozen.tex if some things don't suit your taste
%\include{bibFrozen}

\appendix % switch scrbook parts, numbered: A, B, ...
\begingroup
	\renewcommand*{\cleardoublepage}{\clearpage}
	
\chapter{More information about Foo}


\chapter{More information about Bar}

\section{Detailed List of Evaluation Results}
You get it. % e.g. smaller appendices
	%\include{otherappendixfiles} % for example if you have lengthy tables or the like.
\endgroup

\backmatter % switch scrbook parts, unnumbered
\chapter{Acronyms}
\begin{acronym}[SPARQL]
	% USAGE:
	%\ac{RDF} % first will expand to fullname (short), all subseq. to short
	%\ac*{RDF} % will not mark the acronym as used % note: doesn't work
	%\acf{RDF} % will always expand to fullname (short)
	%\acl{RDF} % will always expand to fullname only
	%\acp{RDF} % makes it plural RDFs...
	%
	% DEFINITION:
	%\acro{ACRO}[short name]{fullname\acroextra{ extra info in this list}}
	% shortname can be used include Fomulas:
	%\acro{chi2}[$\chi^2$]{\texorpdfstring{$\chi^2$-Test}{ChiSquared-Test}}
	
	\acro{CURIE}{Compact URI\acroextra{\footnote{\url{http://www.w3.org/TR/curie/} \lastAccessed}}}
	\acro{LD}{Linked Data\acroextra{\footnote{\url{http://esw.w3.org/topic/LinkedData} \lastAccessed}}}
	\acro{LOD}{Linking Open Data\acroextra{\footnote{\url{http://esw.w3.org/SweoIG/TaskForces/CommunityProjects/LinkingOpenData} \lastAccessed}}}
	\acro{MSE}{Mean Square Error}
	\acro{OWL}{Web Ontology Language\acroextra{\footnote{\url{http://www.w3.org/TR/owl-ref/} \lastAccessed}}}
	\acro{RDF}{Resource Description Framework\acroextra{\footnote{\url{http://www.w3.org/RDF/} \lastAccessed}}}
	\acro{RDFS}{RDF Schema\acroextra{\footnote{\url{http://www.w3.org/TR/rdf-schema/} \lastAccessed}}}
	\acro{SPARQL}{SPARQL Protocol and RDF Query Language\acroextra{\footnote{\url{http://www.w3.org/TR/rdf-sparql-query/} \lastAccessed}}}
	\acro{SKOS}{Simple Knowledge Organization System\acroextra{\footnote{\url{http://www.w3.org/2004/02/skos/} \lastAccessed}}}
	\acro{URI}{Uniform Resource Identifier\acroextra{\footnote{\url{http://tools.ietf.org/html/rfc3986} \lastAccessed}}}
\end{acronym}
 % include a list of acronyms here
\printindex % include the index here

% include the statement about doing the work yourself (if ngerman already you can leave the otherlanguage part)
\begin{otherlanguage}{ngerman}
	\begin{otherlanguage}{ngerman}
% 	Hiermit erkläre ich, dass ich die vorliegende Arbeit selbst angefertigt habe; die aus fremden Quellen direkt oder indirekt übernommenen Gedanken sind als solche kenntlich gemacht.
% 	Die Arbeit wurde bisher keiner Prüfungsbehörde vorgelegt und auch noch nicht veröffentlicht.
	\chapter{Erklärung}

	Hiermit erkläre ich, dass ich die vorliegende Arbeit selbstständig verfasst und keine anderen als die angegebenen Hilfsmittel benutzt habe.
	Die Stellen, die anderen Werken dem Wortlaut oder dem Sinn nach entnommen wurden, habe ich durch die Angabe der Quelle, auch der benutzten Sekundärliteratur, als Entlehnung kenntlich gemacht.
	
	%Teile dieser Arbeit sind während der Bearbeitungszeit bereits in \cite{blabla} veröffentlicht worden.
	\vspace{3em}
	\begin{center}
	\begin{tabular}{cc}
		                                     & \hspace{5cm}    \\
		Kaiserslautern, den \submissiondateDE   & \hrulefill      \\
		                                     & \TheAuthor
	\end{tabular}
	\end{center}
\end{otherlanguage}

\end{otherlanguage}

\end{document}
